% GENERATED FILE. Edit canonical lessons, then run npm run generate:books.
% canonical-source-hash: fnv1a64:f4794f95b41a0f00
% canonical-lessons: GE-C37-dich, GE-C37-ihn, GE-C37-uns, GE-C37-euch, GE-C37-sich, GE-R37-mich-dich-sich

\chapter{The One It Happens To}
\label{ch:ge-objektpronomen}
\hlchaptermodality{50}

\begin{chapteropening}
\textbf{By the end of this chapter:} \emph{I can say who a verb is aimed at with all six object pronouns, and bend a verb back onto its own subject.}

Every verb of address in the book is pointed at every person, and the der-row is run six words deep -- den, einen, keinen, meinen, deinen and now ihn, all moving together. The reflexive costs one word rather than six, because mich, dich, uns and euch already did that job; and freut mich, said since the reader's first introduction as an unexplained block, turns out to have been a reflexive all along.
\end{chapteropening}

\section[dich]{dich — \textquotedblleft{}you,\textquotedblright{} as the one it happens to}

\label{lesson:GE-C37-dich}



\begin{quote}
You own \textbf{mich} (\textquotedblleft{}me\textquotedblright{}) and \textbf{mir} (\textquotedblleft{}to me\textquotedblright{}), and on the \emph{du} side you own \textbf{du}, \textbf{dir} and now \textbf{dein}. One shape is missing from that row.

\begin{itemize}
\raggedright
  \item \emph{Recall:} two chapters back, the joiner that sends the verb to the end — \emph{weil}
\end{itemize}
\end{quote}

\subsection*{You'll want to know: dich}
\begin{quote}
\textbf{dich} — \textquotedblleft{}you\textquotedblright{}, when you are the one being done to.
\end{quote}

Said \emph{dikh}, with the soft ich-sound of \emph{ich}, \emph{mich} and \emph{nicht}.

\begin{quote}
\textbf{Ich sehe dich.} — \textquotedblleft{}I see you.\textquotedblright{}
\end{quote}

\begin{quote}
\textbf{Ich höre dich.} — \textquotedblleft{}I hear you.\textquotedblright{}
\end{quote}

It is \emph{mich}'s partner exactly: \textbf{ich $\to$ mich}, \textbf{du $\to$ dich}. One vowel apart, one job apart, and the pattern is visible without being taught.

\begin{cousinweb}
English once had this whole row and let it collapse. \emph{Thou} was the doer, \emph{thee} was both the done-to and the to-whom, \emph{thine} the possessive. German kept four distinct shapes where English kept three and then lost all of them:

\begin{quote}
\textbf{du} · \textbf{dich} · \textbf{dir} · \textbf{dein}
\end{quote}

Every one begins with the \emph{d} that was English \emph{th}, so the row is predictable from a rule you have had since the second chapter.
\end{cousinweb}

\subsection*{Your turn}
Say these aloud:

\begin{itemize}
\raggedright
  \item \textquotedblleft{}Ich sehe dich.\textquotedblright{} and \textquotedblleft{}Ich höre dich.\textquotedblright{}
  \item \textquotedblleft{}Ich verstehe dich.\textquotedblright{} and \textquotedblleft{}Ich frage dich.\textquotedblright{}
  \item the pair — \textquotedblleft{}Ich sehe dich\textquotedblright{}, \textquotedblleft{}Du siehst mich\textquotedblright{}
  \item the whole \emph{du} row — \textquotedblleft{}du, dich, dir, dein\textquotedblright{}
\end{itemize}

\emph{Twice through:} \textquotedblleft{}Ich sehe dich. Du siehst mich.\textquotedblright{}

\begin{tcolorbox}[breakable,colback=teal!4,colframe=teal!35!black,title={Before you move on}]
Say the four shapes of \emph{du}. (\emph{\textbf{Du, dich, dir, dein}}.) Which one is the done-to? (\emph{\textbf{Dich}}.) Say \textquotedblleft{}I hear you.\textquotedblright{} (\textbf{Ich höre dich.})
\end{tcolorbox}
\section[ihn]{ihn — \textquotedblleft{}him,\textquotedblright{} on the row that always moves}

\label{lesson:GE-C37-ihn}



\begin{quote}
One row of German keeps changing shape and the other two do not. Say \textbf{der $\to$ den}, \textbf{ein $\to$ einen}, \textbf{kein $\to$ keinen}, \textbf{mein $\to$ meinen}. The pronoun does it too.
\end{quote}

\subsection*{You'll want to know: ihn}
\begin{quote}
\textbf{ihn} — \textquotedblleft{}him\textquotedblright{}, and also \textquotedblleft{}it\textquotedblright{} when the \emph{it} is a \emph{der}-word.
\end{quote}

Said \emph{een}: the \textbf{h} is silent and stretches the vowel, as in \emph{ihr} and \emph{Ihnen}.

\begin{quote}
\textbf{Ich sehe ihn.} — \textquotedblleft{}I see him.\textquotedblright{}
\end{quote}

\begin{quote}
\textbf{Der Hund? Ich sehe ihn.} — \textquotedblleft{}The dog? I see it.\textquotedblright{}
\end{quote}

That second sentence is the one an English speaker forgets. A German \emph{der}-word is \emph{ihn}, whatever it is made of, because German picks the pronoun by the article and not by whether the thing is alive.

\textbf{Sie} and \textbf{es} do not move at all — the \emph{die} row and the \emph{das} row stand still here exactly as their articles do.

\begin{cousinweb}
English \textbf{him} looks like a translation of \emph{ihn} and is not the same form. \emph{Him} was originally the \textbf{to-whom} shape, and it swallowed the done-to shape \emph{hine} somewhere in the Middle Ages. German never let that merger happen, which is why it still has \textbf{ihn} for the done-to and \textbf{ihm} for the to-whom.

So the two languages did not lose the same thing. English lost the distinction; German lost nothing and merely stopped being able to explain it briefly.
\end{cousinweb}

\subsection*{Your turn}
Say these aloud:

\begin{itemize}
\raggedright
  \item \textquotedblleft{}Ich sehe ihn.\textquotedblright{} and \textquotedblleft{}Ich höre ihn.\textquotedblright{}
  \item the thing, not the person — \textquotedblleft{}Der Hund? Ich sehe ihn.\textquotedblright{}
  \item \textquotedblleft{}Der Kaffee? Ich trinke ihn.\textquotedblright{} and \textquotedblleft{}Der Wein? Ich sehe ihn.\textquotedblright{}
  \item the row that does not move — \textquotedblleft{}Die Katze? Ich sehe sie.\textquotedblright{}
  \item the three done-to pronouns so far — \textquotedblleft{}mich, dich, ihn\textquotedblright{}
\end{itemize}

\emph{Twice through:} \textquotedblleft{}Ich sehe ihn. Ich sehe sie.\textquotedblright{}

\begin{tcolorbox}[breakable,colback=teal!4,colframe=teal!35!black,title={Before you move on}]
Which pronoun replaces a \emph{der}-word being done to? (\emph{\textbf{Ihn}} — alive or not.) Which two do not change at all? (\emph{\textbf{Sie}} and \emph{\textbf{es}}.) What was English \emph{him} originally? (\textbf{The to-whom form}, which swallowed the done-to one.)
\end{tcolorbox}
\section[uns]{uns — \textquotedblleft{}us,\textquotedblright{} and the one place German is simpler}

\label{lesson:GE-C37-uns}



\begin{quote}
\textbf{Ich} has three shapes — \emph{ich}, \emph{mich}, \emph{mir}. Say them, then take the plural, which has two.
\end{quote}

\subsection*{You'll want to know: uns}
\begin{quote}
\textbf{uns} — \textquotedblleft{}us\textquotedblright{}, and also \textquotedblleft{}to us\textquotedblright{}.
\end{quote}

\begin{quote}
\textbf{Er sieht uns.} — \textquotedblleft{}He sees us.\textquotedblright{}
\end{quote}

\begin{quote}
\textbf{Er gibt uns Kaffee.} — \textquotedblleft{}He gives us coffee.\textquotedblright{}
\end{quote}

One word, both jobs. Where \emph{ich} needed \emph{mich} and \emph{mir}, \textbf{wir} needs only \textbf{uns}, and the language stops distinguishing exactly where it had been insisting.

\begin{cousinweb}
\textbf{uns} and English \textbf{us} are the same word, worn to the same length in both languages. English \emph{us} also does both jobs — \emph{he sees us}, \emph{he gives us coffee} — so this is one row where the two languages agree completely and nothing has to be learned twice.

It is worth noticing \emph{why} the agreement happened. German kept its distinctions in the singular, where the words are short and said constantly; it let them go in the plural, where they are longer and rarer. English did the same thing to \emph{thou} and \emph{thee} four centuries later, only more thoroughly.
\end{cousinweb}

\subsection*{Your turn}
Say these aloud:

\begin{itemize}
\raggedright
  \item \textquotedblleft{}Er sieht uns.\textquotedblright{} and \textquotedblleft{}Er hört uns.\textquotedblright{}
  \item both jobs at once — \textquotedblleft{}Er sieht uns\textquotedblright{} and \textquotedblleft{}Er gibt uns Kaffee\textquotedblright{}
  \item the four done-to pronouns — \textquotedblleft{}mich, dich, ihn, uns\textquotedblright{}
  \item \textquotedblleft{}Verstehst du uns?\textquotedblright{}
\end{itemize}

\emph{Twice through:} \textquotedblleft{}Er sieht uns. Er gibt uns Kaffee.\textquotedblright{}

\begin{tcolorbox}[breakable,colback=teal!4,colframe=teal!35!black,title={Before you move on}]
How many shapes does \emph{wir} have? (\textbf{One} — \emph{uns}, for both jobs.) How many does \emph{ich} have? (\textbf{Three} — \emph{ich, mich, mir}.) Say \textquotedblleft{}he sees us.\textquotedblright{}
\end{tcolorbox}
\section[euch]{euch — \textquotedblleft{}you all,\textquotedblright{} and the word English promoted}

\label{lesson:GE-C37-euch}



\begin{quote}
\textbf{Ihr} is \emph{you all}. It needs a done-to shape, and the shape it has is the most quietly famous word in this chapter.
\end{quote}

\subsection*{You'll want to know: euch}
\begin{quote}
\textbf{euch} — \textquotedblleft{}you all\textquotedblright{}, when you are the ones being done to. It also covers \emph{to you all}, exactly as \emph{uns} covers both jobs.
\end{quote}

Said \emph{oykh}: the \textbf{eu} is the \emph{oy} of English \emph{boy}, and the \textbf{ch} is the soft one.

\begin{quote}
\textbf{Ich sehe euch.} — \textquotedblleft{}I see you all.\textquotedblright{}
\end{quote}

\begin{cousinweb}
English had this word. Old English \textbf{ēow} was the done-to shape of \textbf{ġē} (\textquotedblleft{}ye\textquotedblright{}), and \emph{ēow} is exactly what German's \emph{euch} is.

Then English did something strange. It let \emph{ēow} — the \textbf{object} form — take over the subject job as well, so \emph{ye} died out and \textbf{you} became the only word left for the second person, singular or plural, doer or done-to. Modern English addresses everybody with what was once a plural object pronoun.

German still has the whole set standing: \textbf{du, dich, dir} for one person, \textbf{ihr, euch} for several, \textbf{Sie, Ihnen} for a stranger. English boiled all of that down to \textbf{you}, and \emph{y'all} is the attempt to build the plural back.
\end{cousinweb}

\subsection*{Your turn}
Say these aloud:

\begin{itemize}
\raggedright
  \item \textquotedblleft{}Ich sehe euch.\textquotedblright{} and \textquotedblleft{}Ich höre euch.\textquotedblright{}
  \item the five done-to pronouns — \textquotedblleft{}mich, dich, ihn, uns, euch\textquotedblright{}
  \item the three ways to address people — \textquotedblleft{}dich, euch, Ihnen\textquotedblright{}
  \item the English word all of them collapsed into — \emph{you}
\end{itemize}

\emph{Twice through:} \textquotedblleft{}mich, dich, ihn, uns, euch.\textquotedblright{}

\begin{tcolorbox}[breakable,colback=teal!4,colframe=teal!35!black,title={Before you move on}]
What is the done-to shape of \emph{ihr}? (\emph{\textbf{Euch}}.) Which English word is it? (\emph{\textbf{You}} — by way of Old English \emph{ēow}.) Which English word did \emph{ēow} push out? (\emph{\textbf{Ye}}.) Say all five done-to pronouns you now own.
\end{tcolorbox}
\section[sich]{sich — the verb bent back on itself}

\label{lesson:GE-C37-sich}



\begin{quote}
You have said \textbf{freut mich} since the second chapter without asking what the \emph{mich} is doing there. This lesson answers it.
\end{quote}

\begin{grammarlens}[title={Grammar Lens: when the doer and the done-to are the same}]
\begin{quote}
When the person doing something is also the person it happens to, German puts the done-to pronoun in anyway.
\end{quote}

\begin{quote}
\textbf{Ich wasche mich.} — \textquotedblleft{}I wash.\textquotedblright{} Literally \emph{I wash me}.
\end{quote}

English drops the pronoun in most of these — \emph{I wash}, \emph{I shave}, \emph{I dress} — and German never does. Every done-to pronoun you own is already the right one: \emph{mich}, \emph{dich}, \emph{uns}, \emph{euch}. Nothing new is needed for four of the six.
\end{grammarlens}

\subsection*{You'll want to know: sich}
\begin{quote}
\textbf{sich} — the one shape German uses for \emph{himself}, \emph{herself}, \emph{itself}, \emph{themselves} and formal \emph{yourself}.
\end{quote}

Said \emph{zikh}, with a \emph{z} at the front and the soft ich-sound at the back.

\begin{quote}
\textbf{Er wäscht sich.} · \textbf{Sie wäscht sich.} · \textbf{Sie waschen sich.}
\end{quote}

Where English needs four different words with \emph{-self} welded on, German needs one, and it never changes.

\begin{cousinweb}
\textbf{sich} is Latin \textbf{sē} and French \textbf{se}, all descended from Proto-Indo-European \emph{\textbf{*swe-}}, \textquotedblleft{}one's own\textquotedblright{}. English lost the pronoun and kept the root inside words built on it: \textbf{self}, \textbf{suicide}, and — by a long road — \textbf{sibling}, a person of one's own kin.

Which finally explains \textbf{freut mich}: \emph{freuen} is one of these verbs, so \emph{es freut mich} is \textquotedblleft{}it gladdens me\textquotedblright{}, and the greeting you have said a hundred times was a reflexive construction all along.
\end{cousinweb}

\subsection*{Your turn}
Say these aloud:

\begin{itemize}
\raggedright
  \item \textquotedblleft{}Ich wasche mich.\textquotedblright{} and \textquotedblleft{}Er wäscht sich.\textquotedblright{}
  \item the whole set, four owned and one new — \textquotedblleft{}mich, dich, sich, uns, euch\textquotedblright{}
  \item the greeting, taken apart — \textquotedblleft{}freut mich\textquotedblright{}
  \item \textquotedblleft{}Sie freut sich.\textquotedblright{}
\end{itemize}

\emph{Twice through:} \textquotedblleft{}Ich wasche mich. Er wäscht sich.\textquotedblright{}

\begin{tcolorbox}[breakable,colback=teal!4,colframe=teal!35!black,title={Before you move on}]
How many English words does \emph{sich} cover? (\textbf{Four or more} — himself, herself, itself, themselves.) Which four done-to pronouns did you already own for this job? (\emph{\textbf{Mich, dich, uns, euch}}.) And what has \emph{freut mich} been all along? (\textbf{A reflexive} — \textquotedblleft{}it gladdens me\textquotedblright{}.)
\end{tcolorbox}
\section[Practice]{mich, dich, sich — the chapter, run cold}

\label{lesson:GE-R37-mich-dich-sich}



\begin{quote}
Say the six done-to pronouns. Then name the two that did not have to be learned as new words, and say why.
\end{quote}

\subsection*{Your turn: point every verb at somebody}
Six verbs of address, six pronouns. Run the grid by ear, not by reading it.

Say these aloud:

\begin{itemize}
\raggedright
  \item \textquotedblleft{}Ich sehe dich. Ich sehe ihn. Ich sehe euch.\textquotedblright{}
  \item \textquotedblleft{}Ich höre dich. Er hört uns. Wir hören euch.\textquotedblright{}
  \item \textquotedblleft{}Ich verstehe dich.\textquotedblright{} and \textquotedblleft{}Verstehst du mich?\textquotedblright{}
  \item \textquotedblleft{}Ich frage dich.\textquotedblright{} and \textquotedblleft{}Ich helfe dir.\textquotedblright{} — one takes the done-to shape, one takes the to-whom
  \item the two verbs of the page — \textquotedblleft{}Ich lese es. Ich schreibe es.\textquotedblright{}
\end{itemize}

\subsection*{Your turn: the distant band — the row that moves}
Six words now move together on one row. Say the article, then the whole family, then the pronoun that replaces it.

Say these aloud:

\begin{itemize}
\raggedright
  \item \textquotedblleft{}der Hund $\to$ den Hund, einen Hund, keinen Hund, meinen Hund $\to$ ihn\textquotedblright{}
  \item \textquotedblleft{}der Kaffee $\to$ den Kaffee $\to$ ihn\textquotedblright{}; \textquotedblleft{}der Tee $\to$ den Tee $\to$ ihn\textquotedblright{}; \textquotedblleft{}der Wein $\to$ den Wein $\to$ ihn\textquotedblright{}
  \item the people — \textquotedblleft{}den Bruder, den Vater, den Freund $\to$ ihn\textquotedblright{}
  \item the body — \textquotedblleft{}den Arm, den Fuß, den Finger, den Kopf, den Mund $\to$ ihn\textquotedblright{}
  \item the two rows that stand still — \textquotedblleft{}die Katze $\to$ sie\textquotedblright{}, \textquotedblleft{}die Milch $\to$ sie\textquotedblright{}, \textquotedblleft{}das Brot $\to$ es\textquotedblright{}
\end{itemize}

\subsection*{Your turn: bent back}
Say these aloud:

\begin{itemize}
\raggedright
  \item \textquotedblleft{}Ich wasche mich. Du wäschst dich. Er wäscht sich.\textquotedblright{}
  \item \textquotedblleft{}Wir waschen uns. Ihr wascht euch.\textquotedblright{}
  \item the greeting that turns out to be one of these — \textquotedblleft{}Freut mich.\textquotedblright{}
\end{itemize}

\begin{tcolorbox}[breakable,colback=teal!4,colframe=teal!35!black,title={Before you move on}]
Which pronoun replaces a \emph{der}-word being done to? (\emph{\textbf{Ihn}}.) Which two pronouns never change shape? (\emph{\textbf{Sie}} and \emph{\textbf{es}}.) How many English words does \emph{sich} cover? (\textbf{Four or more}.) And which English pronoun is \emph{euch}? (\emph{\textbf{You}} — the object form that took the whole job.)
\end{tcolorbox}
